\Ch{Classes}{Class}

\Sec{Member Functions}{Class.MemberFunctions}

\Sub{Non-static Member Functions}{Class.MemberFunctions.NonStatic}

% Partial subclause text introduced in proposal 0015.

\p A non-static member function may be declared \texttt{const}, such a function
is called a \textit{const member function}. The \texttt{const} qualifier of a
const member function applies to the \texttt{this} reference, and affects the
type of the member function.

\Sec{Conversion Functions}{Class.MemberFunctions.Conversion}

\begin{grammar}
  \define{conversion-function-id}\br
  \terminal{operator} conversion-type-id\br
\end{grammar}

\p A member function of a class named as a \textit{conversion-function-id}
specifies a conversion from the class type to the type specified by the
\textit{conversion-type-id}. Such functions are called \textit{conversion
functions}. No return type is specified. The type of the function is a function
taking no parameters returning \textit{conversion-type-id}.

\p A conversion function may never be used to convert an object to the same
type, to a base class of the type, or to void.

\p The \texttt{explicit} keyword may be applied to a conversion function
declaration (\ref{Decl.Spec.Fct}). Such functions are said to be
\textit{explicit conversion functions}, which are only considered as a
user-defined conversion for direct-initialization or explicit cast expressions.
